\documentclass[11pt]{article}
% ===========================================================================
%  Shared preamble for the AP Calculus AB exam-preparation worksheets.
%  Input right after \documentclass. Files beginning with "_" are not compiled
%  on their own. Navy and gold, to match the course's lesson pages.
% ===========================================================================
\usepackage[letterpaper,top=2.3cm,bottom=2.2cm,left=1.7cm,right=1.7cm,headheight=28pt,footskip=30pt]{geometry}
\usepackage{amsmath,amssymb}
\usepackage[T1]{fontenc}
\usepackage[utf8]{inputenc}
\usepackage{lmodern}
\usepackage{xcolor}
\usepackage[most]{tcolorbox}
\usepackage{enumitem}
\usepackage{fancyhdr}
\usepackage{array}
\usepackage{tikz}
\usepackage{pgfplots}
\pgfplotsset{compat=1.18}

\definecolor{navy}{HTML}{161E54}
\definecolor{ink}{HTML}{27357E}
\definecolor{gold}{HTML}{B98A00}
\definecolor{goldbg}{HTML}{FFFDF3}
\definecolor{inkbg}{HTML}{F6F7FD}
\definecolor{rule}{HTML}{C9CFF0}
\definecolor{okgreen}{HTML}{1B7A44}
\definecolor{okbg}{HTML}{F1F8F3}

\pagestyle{fancy}
\fancyhf{}
\fancyhead[L]{\footnotesize Name: \underline{\hspace{4cm}}}
\fancyhead[C]{\textbf{\color{navy}Integration Academy}\\[-2pt]{\scriptsize Calculus AB \textperiodcentered{} AP\textsuperscript{\textregistered} Exam Preparation}}
\fancyfoot[L]{\tiny\color{black!55}\parbox[t]{0.78\textwidth}{AP\textsuperscript{\textregistered} is a registered trademark of the College Board, which was not involved in the production of, and does not endorse, this material. \textcopyright{} Integration Academy.}}
\fancyfoot[R]{\scriptsize\thepage}
\renewcommand{\headrulewidth}{1.2pt}
\renewcommand{\headrule}{\hbox to\headwidth{\color{gold}\leaders\hrule height \headrulewidth\hfill}}
\renewcommand{\footrulewidth}{0pt}
\setlength{\parindent}{0pt}
\setlength{\parskip}{4pt}

% Title block: \wstitle{code}{title}{unit line}
\newcommand{\wstitle}[3]{\begin{center}{\small\bfseries\color{gold}WORKSHEET #1}\\[3pt]{\LARGE\bfseries\color{navy}#2}\\[3pt]{\small\color{black!60}#3}\end{center}\vspace{2pt}}
% Section heading
\newcommand{\sech}[1]{\par\vspace{10pt}{\large\bfseries\color{navy}#1}\par\nobreak\vspace{2pt}{\color{rule}\hrule height0.8pt}\vspace{7pt}}
% A level tag printed beside a group of questions
\newcommand{\level}[1]{\par\vspace{6pt}{\footnotesize\bfseries\color{gold}\MakeUppercase{#1}}\par\nobreak\vspace{2pt}}

\newtcolorbox{ideabox}{colframe=ink,colback=inkbg,coltitle=white,colbacktitle=ink,fonttitle=\bfseries\small,title={Key ideas},arc=2pt,boxrule=0.8pt,left=8pt,right=8pt}
\newtcolorbox{example}[1]{colframe=rule,colback=white,coltitle=navy,colbacktitle=inkbg,fonttitle=\bfseries\small,title={#1},arc=2pt,boxrule=0.8pt,borderline west={3pt}{0pt}{ink},left=10pt,right=8pt}
\newtcolorbox{frbox}[1]{colframe=gold,colback=goldbg,coltitle=white,colbacktitle=gold,fonttitle=\bfseries\small,title={#1},arc=2pt,boxrule=0.8pt,breakable,left=8pt,right=8pt}
\newtcolorbox{solbox}{colframe=okgreen,colback=okbg,coltitle=white,colbacktitle=okgreen,fonttitle=\bfseries\small,title={Worked solutions},arc=2pt,boxrule=0.8pt,breakable,left=8pt,right=8pt}
\newtcolorbox{scorebox}[1]{colframe=navy,colback=white,coltitle=white,colbacktitle=navy,fonttitle=\bfseries\small,title={#1},arc=2pt,boxrule=0.8pt,breakable,left=8pt,right=8pt}

\newcommand{\soln}{\par\smallskip\textbf{\color{ink}Solution.}\ }
\newcommand{\result}[1]{\par\smallskip\textbf{#1}}
% blank working space of a given height
\newcommand{\work}[1]{\par\vspace{#1}}
% a point value in the margin of a scoring guide
\newcommand{\pts}[1]{\hfill{\footnotesize\bfseries\color{gold}[#1]}}

% Coordinate grid for graph questions: \begin{apgraph}{xmin}{xmax}{ymin}{ymax} ... \end{apgraph}
\newenvironment{apgraph}[4]{\begin{center}\begin{tikzpicture}\begin{axis}[
  axis lines=middle, xmin=#1, xmax=#2, ymin=#3, ymax=#4,
  grid=both, grid style={black!12}, xtick distance=1, ytick distance=1,
  ticklabel style={font=\scriptsize}, xlabel={$x$}, ylabel={$y$},
  width=10cm, height=7.2cm, samples=80, clip=false,
  every axis plot/.append style={line width=1.1pt, color=ink}]}{\end{axis}\end{tikzpicture}\end{center}}
\newcommand{\closeddot}[2]{\addplot[only marks, mark=*, mark size=2.2pt, color=ink] coordinates {(#1,#2)};}
\newcommand{\opendot}[2]{\addplot[only marks, mark=*, mark size=2.2pt, color=ink, mark options={fill=white}] coordinates {(#1,#2)};}

% Multiple-choice options: four across, or stacked when the options are long.
\newcommand{\opts}[4]{\par\vspace{3pt}\begin{tabular}{@{}p{0.235\linewidth}p{0.235\linewidth}p{0.235\linewidth}p{0.235\linewidth}@{}}(A)\ #1 & (B)\ #2 & (C)\ #3 & (D)\ #4\end{tabular}\par\vspace{9pt}}
\newcommand{\optsv}[4]{\par\vspace{3pt}\begin{tabular}{@{}l@{\ \ }p{0.9\linewidth}@{}}(A) & #1\\(B) & #2\\(C) & #3\\(D) & #4\end{tabular}\par\vspace{9pt}}

\fancyhead[R]{\footnotesize\color{black!60}Worksheet 1.4}
\begin{document}

\wstitle{1.4}{Limit Laws and Direct Substitution}{Unit 1: Limits and Continuity}

\begin{ideabox}
\begin{itemize}[leftmargin=*,itemsep=2pt,topsep=2pt]
  \item If $\displaystyle\lim_{x\to c}f(x)=L$ and $\displaystyle\lim_{x\to c}g(x)=M$ are real numbers, limits of sums, differences, products and constant multiples are found termwise.
  \item The quotient law requires $M\neq0$.
  \item For polynomial, rational, root, trigonometric, exponential and logarithmic functions, $\displaystyle\lim_{x\to c}f(x)=f(c)$ at every $c$ in the domain.
  \item $\displaystyle\lim_{x\to c}f\big(g(x)\big)=f\big(\displaystyle\lim_{x\to c}g(x)\big)$ only when $f$ is continuous at that inner limit.
\end{itemize}
\end{ideabox}

\sech{Worked examples}
\begin{example}{Example 1 \textperiodcentered{} Direct substitution}
Evaluate $\displaystyle\lim_{x\to -1}\dfrac{2x^3-x+4}{x^2+3}$.\soln The denominator at $x=-1$ is $4\neq0$, so substitute: $\dfrac{-2+1+4}{4}$.\result{The limit is $\dfrac34$.}
\end{example}

\begin{example}{Example 2 \textperiodcentered{} Laws with unknown functions}
Given $\displaystyle\lim_{x\to 1}f(x)=3$ and $\displaystyle\lim_{x\to 1}g(x)=-4$, find $\displaystyle\lim_{x\to 1}\dfrac{[f(x)]^2-2g(x)}{f(x)+g(x)+5}$.\soln Numerator: $9+8=17$. Denominator: $3-4+5=4\neq0$.\result{The limit is $\dfrac{17}{4}$.}
\end{example}

\begin{example}{Example 3 \textperiodcentered{} Working backwards}
Suppose $\displaystyle\lim_{x\to 5}\dfrac{f(x)-2}{x-5}=7$. Find $\displaystyle\lim_{x\to 5}f(x)$.\soln $f(x)-2=\dfrac{f(x)-2}{x-5}\cdot(x-5)\to7\cdot0=0$.\result{$\displaystyle\lim_{x\to 5}f(x)=2$.}
\end{example}

\newpage
\sech{Practice}

\level{Core}
\begin{enumerate}[leftmargin=*,itemsep=2pt]
  \item Evaluate $\displaystyle\lim_{x\to 2}\left(x^4-3x\right)$.\work{2.3cm}
  \item Evaluate $\displaystyle\lim_{x\to \pi/4}\sin x\cos x$.\work{2.3cm}
  \item Evaluate $\displaystyle\lim_{x\to 0}\dfrac{3e^x-2}{x+1}$.\work{2.3cm}
  \item Evaluate $\displaystyle\lim_{x\to 4}\sqrt{x^2+9}$.\work{2.3cm}
\end{enumerate}

\level{Exam level}
\begin{enumerate}[leftmargin=*,itemsep=2pt,start=5]
  \item Given $\displaystyle\lim_{x\to 3}f(x)=-2$ and $\displaystyle\lim_{x\to 3}g(x)=6$, find $\displaystyle\lim_{x\to 3}\big[g(x)-(f(x))^2\big]f(x)$.\work{2.8cm}
  \item With the same $f$ and $g$, find $\displaystyle\lim_{x\to 3}\sqrt{g(x)+f(x)+5}$.\work{2.8cm}
  \item With the same $f$ and $g$, what can be said about $\displaystyle\lim_{x\to 3}\dfrac{f(x)}{g(x)-6}$? Explain.\work{2.8cm}
  \item Suppose $\displaystyle\lim_{x\to -1}\dfrac{f(x)+4}{x+1}=3$. Find $\displaystyle\lim_{x\to -1}f(x)$.\work{2.8cm}
\end{enumerate}

\begin{frbox}{Question 9 \textperiodcentered{} Free-response practice}
Let $f(x)=4+(x-2)^2$ and $g(u)=\begin{cases}u^2,&u\le4\\10-u,&u>4.\end{cases}$
\begin{enumerate}[label=(\alph*),leftmargin=*,itemsep=2pt]
  \item Find $\displaystyle\lim_{x\to 2}f(x)$, and state whether $f(x)$ approaches this value from above, from below, or from both sides.\work{2cm}
  \item Explain why $\displaystyle\lim_{x\to 2}g\big(f(x)\big)$ cannot be found by computing $g\big(\displaystyle\lim_{x\to 2}f(x)\big)$.\work{2cm}
  \item Find $\displaystyle\lim_{x\to 2}g\big(f(x)\big)$.\work{2cm}
  \item Find $\displaystyle\lim_{x\to 2}g\big(k(x)\big)$ where $k(x)=4-(x-2)^2$.\work{2cm}
\end{enumerate}
\end{frbox}

\newpage
\begin{solbox}
\begin{enumerate}[leftmargin=*,itemsep=4pt]
  \item $16-6=10$.
  \item $\dfrac{\sqrt2}{2}\cdot\dfrac{\sqrt2}{2}=\dfrac12$.
  \item $\dfrac{3-2}{1}=1$.
  \item $\sqrt{25}=5$.
  \item $(6-4)(-2)=-4$.
  \item $\sqrt{6-2+5}=3$.
  \item It does not exist. The denominator tends to 0 while the numerator tends to $-2\neq0$, so the quotient is unbounded; the quotient law does not apply.
  \item $f(x)+4=\dfrac{f(x)+4}{x+1}(x+1)\to3\cdot0=0$, so the limit is $-4$.
  \item \begin{enumerate}[label=(\alph*),leftmargin=*,itemsep=3pt]
    \item $4$. Since $(x-2)^2>0$ for $x\neq2$, $f(x)>4$: it approaches 4 from above only.
    \item $g$ is not continuous at $u=4$: $\displaystyle\lim_{u\to 4^-}g(u)=16$ but $\displaystyle\lim_{u\to 4^+}g(u)=6$.
    \item $f(x)\to4^+$, so the limit is $\displaystyle\lim_{u\to 4^+}g(u)=6$.
    \item Now $k(x)<4$ for $x\neq2$, so $k(x)\to4^-$ and the limit is $16$.
  \end{enumerate}
\end{enumerate}
\end{solbox}

\end{document}
